\documentclass[11pt,a4paper]{article}

% Packages
\usepackage[utf8]{inputenc}
\usepackage[T1]{fontenc}
\usepackage{lmodern}
\usepackage[margin=1in]{geometry}
\usepackage{graphicx}
\usepackage{hyperref}
\usepackage{xcolor}
\usepackage{listings}
\usepackage{booktabs}
\usepackage{tabularx}
\usepackage{longtable}
\usepackage{amsmath}
\usepackage{amssymb}
\usepackage{bytefield}
\usepackage{fancyhdr}
\usepackage{titlesec}
\usepackage{enumitem}
\usepackage{tcolorbox}
\usepackage{float}
\usepackage{caption}
\usepackage{subcaption}
\usepackage{multicol}
\usepackage{tikz}
\usetikzlibrary{shapes,arrows,positioning,calc}

% Color definitions
\definecolor{skyblue}{RGB}{0,102,204}
\definecolor{darkblue}{RGB}{0,51,102}
\definecolor{codebg}{RGB}{245,245,245}
\definecolor{codegreen}{RGB}{0,128,0}
\definecolor{codegray}{RGB}{128,128,128}
\definecolor{codepurple}{RGB}{128,0,128}
\definecolor{warnorange}{RGB}{255,140,0}
\definecolor{alertred}{RGB}{204,0,0}

% Hyperref setup
\hypersetup{
	colorlinks=true,
	linkcolor=darkblue,
	urlcolor=skyblue,
	citecolor=darkblue,
	pdftitle={Kestrel Protocol -- Compliance and Regulatory Engineering Report},
	pdfauthor={Krishnashis Das},
}

% Code listing style
\lstdefinestyle{cstyle}{
	backgroundcolor=\color{codebg},
	basicstyle=\ttfamily\footnotesize,
	breaklines=true,
	captionpos=b,
	commentstyle=\color{codegreen},
	keywordstyle=\color{skyblue}\bfseries,
	stringstyle=\color{codepurple},
	numberstyle=\tiny\color{codegray},
	numbers=left,
	numbersep=5pt,
	frame=single,
	rulecolor=\color{codegray},
	language=C,
	showstringspaces=false,
	tabsize=4,
	morekeywords={uint8_t,uint16_t,uint32_t,int8_t,int16_t,int32_t,bool,size_t,ssize_t,true,false},
}
\lstset{style=cstyle}

% Header/Footer
\setlength{\headheight}{14pt}
\pagestyle{fancy}
\fancyhf{}
\fancyhead[L]{\textcolor{darkblue}{\textbf{Kestrel Compliance Report}}}
\fancyhead[R]{\nouppercase{\leftmark}}
\fancyfoot[L]{\small © Wingspann Global Pvt Ltd}
\fancyfoot[R]{\thepage}
\renewcommand{\headrulewidth}{0.4pt}
\renewcommand{\footrulewidth}{0.4pt}
\renewcommand{\sectionmark}[1]{\markboth{#1}{}}

% tcolorbox environments
\newtcolorbox{notebox}{colback=blue!5!white,colframe=skyblue,title=\textbf{Note},fonttitle=\bfseries}
\newtcolorbox{warnbox}{colback=orange!5!white,colframe=warnorange,title=\textbf{Warning},fonttitle=\bfseries}
\newtcolorbox{alertbox}{colback=red!5!white,colframe=alertred,title=\textbf{Security},fonttitle=\bfseries}

% Title formatting
\titleformat{\section}{\Large\bfseries\color{darkblue}}{\thesection}{1em}{}
\titleformat{\subsection}{\large\bfseries\color{skyblue}}{\thesubsection}{1em}{}
\titleformat{\subsubsection}{\normalsize\bfseries\color{darkblue}}{\thesubsubsection}{1em}{}

\begin{document}

% ============================================================
% TITLE PAGE
% ============================================================
\begin{titlepage}
	\centering
	{\includegraphics[width=0.8\textwidth]{Icon.png}\par}
	\vspace*{2cm}
	{\Huge\bfseries\textcolor{darkblue}{KESTREL}\par}
	\vspace{0.5cm}
	{\LARGE\textcolor{skyblue}{}\par}
	\vspace{0.3cm}
	{\large\textcolor{codegray}{Compliance and Regulatory Engineering Report}\par}
	\vspace{2cm}
	{\large Aerospace Compliance Architecture, Implementations, and Validations\par}
	\vspace{3cm}
	
	{\normalsize
		\textbf{Theorized and Developed by:} Krishnashis Das\\
		GitHub: \href{https://github.com/I-am-Krish}{\texttt{I-am-Krish}}\\[1em]
		\textbf{Co-Developed by:} Aadarsh Ajay Kumar Sinha\\
		GitHub: \href{https://github.com/Monarch666}{\texttt{Monarch666}}\\[1em]
		\textbf{Document Revision:}\\
		\vspace{0.8cm}
		\LARGE\textcolor{alertred}{\textbf{Strictly Confidential}}
	}
	
	\vfill
	{\footnotesize\textcolor{codegray}{This document details the engineering modifications and implementation strategies executed to achieve international aerospace regulatory compliance (DO-362A, DO-377A, DGCA NPNT, ASTM F3411, JARUS SORA OSO\#06).}}
\end{titlepage}

% ============================================================
% DOCUMENT CONTROL
% ============================================================
\newpage
\section*{Document Control}
\addcontentsline{toc}{section}{Document Control}

\begin{table}[H]
	\renewcommand{\arraystretch}{1.5}
	\begin{tabularx}{\textwidth}{@{}lX@{}}
		\toprule
		\textbf{Field} & \textbf{Value} \\
		\midrule
		\textbf{Document Title} & Kestrel Protocol Compliance Engineering Report \\
		\textbf{Document ID} & UL-COMP-002 \\
		%\textbf{Version} & \\
		%\textbf{Status} & Multi-Variant Release / Audit Closed \\
		\midrule
		\textbf{Author} & Krishnashis Das \\
		\textbf{Co-Author} & Aadarsh Ajay Kumar Sinha \\
		\textbf{Organization} & Wingspann Global Pvt Ltd \\
		\midrule
		\textbf{Release Date} & April 15, 2026 \\
		\midrule
		\textbf{Distribution} & Internal Use Only \\
		\textbf{Classification} & \textcolor{alertred}{\textbf{Strictly Confidential}} \\
		\bottomrule
	\end{tabularx}
\end{table}

\vspace{1cm}

\begin{alertbox}
	\textbf{Document Security Notice:} This analytical compliance report details strict architectural maneuvers involving cryptographic gates and internal payload strategies. Unauthorized distribution or exploitation is strongly prohibited.
\end{alertbox}

\newpage
\tableofcontents
\newpage

% ============================================================
% 1. PROTOCOL PHILOSOPHY & STRUCTURAL APPROACH
% ============================================================
\section{Philosophy and Structural Approach to Compliance}

The fundamental challenge in modernizing Kestrel from a proprietary, hermetically sealed encrypted tunnel into an internationally compliant aerospace framework was resolving the inherent contradictions between strict military cryptography and open-access regulatory tracking.

To overcome this, we employed an \textbf{Additive Shim Layer} philosophy.

\subsection{The Additive Shim Layer Mechanism}
Instead of stripping down or complicating the hyper-optimized foundational core, every compliance feature was engineered as a decoupled modular loop.
\begin{itemize}
	\item \textbf{Preservation of Core Metrics:} The primary zero-copy parsing engine, the strictly deterministic O(1) memory pool allocation, and the ChaCha20-Poly1305 AEAD symmetric cryptography pipelines were deliberately insulated from these revisions.
	\item \textbf{Uncompromised Performance:} Real-time telemetry parse speeds (125$\mu$s) and overall end-to-end latencies stay securely under sub-millisecond ranges without friction from newly introduced tracking variables.
	\item \textbf{State Machine Overlays:} We added specific parser interceptions in the simulators rather than mutating core state tracking. This ensures Kestrel maintains true bidirectional compliance seamlessly across drastically different domains (e.g. NATO MIL-STANAG networks vs. civilian ASTM beacons).
\end{itemize}

% ============================================================
% 2. STATE OVERVIEW MATRIX
% ============================================================
\section{Protocol Compliance State Matrix}

The following matrix provides a quick-glance visualization of Kestrel's current compliance limits, detailing the targeted domains and the status of structural implementations.

\begin{table}[H]
	\centering
	\renewcommand{\arraystretch}{1.4}
	\begin{tabularx}{\textwidth}{@{} | l | l | X | c |@{}}
		\toprule
		\textbf{Standard} & \textbf{Domain Focus} & \textbf{Architectural Modification \& Status} & \textbf{Completion} \\
		\midrule
		\textbf{RTCA DO-362A} & Terrestrial C2 & Configurable action/timeout, GCS authority loop, latency gate & \textbf{$\sim$85\%} \\
		\hline
		\textbf{RTCA DO-377A} & BLOS SATCOM & 4-slot window, exponential backoff, RTT budget script & \textbf{$\sim$80\%} \\
		\hline
		\textbf{DGCA NPNT} & Authenticated Flight & Ed25519 PA verification, geofence, time window, live test keys & \textbf{$\sim$75\%} \\
		\hline
		\textbf{ASTM F3411} & Open Public ID & Isolated RID module, ENCRYPT\_NEVER policy, 1 Hz beacon & \textbf{$\sim$80\%} \\
		\hline
		\textbf{STANAG 4586} & Interoperability & Identifier hooks, stream-type partitioning, LOI scaffolding & \textbf{$\sim$30\%} \\
		\hline
		\textbf{STANAG 4609} & Motion Imagery & ST0601 KLV metadata generation, MPEG-TS multiplexing & \textbf{$\sim$85\%} \\
		\hline
		\textbf{JARUS SORA OSO\#06} & C2 Cyber Security & 64-entry security event ring buffer (\texttt{kestrel\_sora.c}); hooks on replay/MAC/auth/key events; \texttt{ks\_sora\_is\_compliant()} assertion; audit dump on shutdown & \textbf{$\sim$95\%} \\
		%\textbf{NASA-STD-1006} & Space System Protection & SSPR-1 inherently met (ChaCha20-Poly1305); SSPR-2/3/4 not implemented & \textbf{$\sim$25\%} \\
		\bottomrule
	\end{tabularx}
\end{table}

% ============================================================
% 3. DETAILED EXAMINATIONS
% ============================================================
\section{Detailed Compliance Specifications}

This section breaks down exactly what compliance each standard dictates, its relevance to Kestrel, the precise surgical code adjustments executed alongside clarifying logic snippets, and the concluding resulting behavior.

\subsection{RTCA DO-362A: Terrestrial C2 Data Link Failsafes}

\subsubsection{What It States}
RTCA DO-362A strictly regulates Unmanned Aircraft System (UAS) Command and Control (C2) terrestrial limits. Specifically, it explicitly \textbf{prohibits hardcoded autonomic failsafe responses}. In the event of a C2 link loss, an unprompted "Return-to-Launch" could inadvertently route the UAS through active civilian airspace. The governing authority (GCS) must definitively parameterize and update the exact Lost-Link geometry (Land, Hover, RTL) and timeout threshold dynamically.

\subsubsection{Relevance to Kestrel}
Earlier versions of Kestrel possessed a rigid fail-safe fallback mechanism encoded natively into the simulation loop (\texttt{uav\_simulator.c}). If the C2 heartbeat perished for a static 3000ms, the system unilaterally initiated an RTL. This directly violated DO-362A's dynamic authority requirements.

\subsubsection{Implemented Code Changes}

We extended \texttt{ks\_heartbeat\_t} payload to carry dynamic variables. The GCS is programmed to actively transmit these down effectively overriding the UAV variables in real-time.

\begin{lstlisting}[language=C, caption={UAV Heartbeat Parser State Lock (\texttt{uav\_simulator.c})}]
// Inside uav_simulator.c packet parser switch
case KS_MSG_HEARTBEAT:
    ks_deserialize_heartbeat(&hb, &pkt.payload[0]);
    // Dynamically lock UAV failsafe limits to live GCS dictates
    g_failsafe_action = hb.lost_link_action;
    g_failsafe_timeout_s = hb.lost_link_timeout_s;
    break;
\end{lstlisting}

\subsubsection{Intended Protocol Behavior and Impact}
The UAV now dynamically adapts to the commanding station. A GCS enforces its failsafe conditions instantly onto the UAV. If the GCS goes completely dead, the UAV acts upon the very last known verifiable instruction set given by the central authority. It guarantees DO-362A compliant C2 authority looping, securing deployment in contested airways.

\vspace{0.5cm}

\subsection{RTCA DO-377A: Beyond Line-of-Sight (BLOS) Satellite Latency}

\subsubsection{What It States}
DO-377A sets boundaries for BLOS networking, specifically when satellite communication induces drastic propagation delays. Round trip margins can frequently stretch beyond 1.5 seconds. Systems must be hardened against high-latency network drops without bottlenecking commands.

\subsubsection{Relevance to Kestrel}
Prior versions relied on single-instance locking—the Ground Station transmitted one command, locked its pipeline, and idled pending the \texttt{CMD\_ACK}. For terrestrial lines, this is unnoticeable. On LEO/GEO satellite uplinks, it creates a fatal head-of-line blocking cascade.

\subsubsection{Implemented Code Changes}

We decoupled the transmission lock and built an asynchronous sliding window.

\begin{lstlisting}[language=C, caption={Sliding Window Allocation Structure (\texttt{gcs\_receiver.c})}]
// Inside gcs_receiver.c pipeline management
#define DO_377A_MAX_WINDOW_SIZE 4 

typedef struct {
    bool active;
    uint16_t msg_id;
    uint32_t transmit_time_ms;
    uint32_t retry_delay_ms; // Escalates exponentially
    uint8_t retries;         // Capped at DO_377A_MAX_RETRIES
} ks_cmd_slot_t;
\end{lstlisting}

\subsubsection{Intended Protocol Behavior and Impact}
Operating Kestrel under extreme delays is now remarkably fluid. Operators deploying concurrent flight vectors no longer experience software lock-ups, establishing robust parallel transmission layers directly compliant with DO-377A high-latency SATCOM models while preserving timeout resilience.

\vspace{0.5cm}

\subsection{DGCA NPNT: Digital Airspace Clearance Gates}

\subsubsection{What It States}
India’s Directorate General of Civil Aviation requires \textbf{No-Permission, No-Takeoff} (NPNT) functionality. The system must intrinsically deny physical arming attempts natively within its controller unless a cryptographically verifiable Permission Artifact (PA)—certified by the central authority—is verified.

\subsubsection{Relevance to Kestrel}
Kestrel lacked structural procedural hardware barriers outside its symmetric AEAD session streams.

\subsubsection{Implemented Code Changes}

We added an intercept handler directly inside the telemetry validation pipeline preventing internal subsystem bridging if the crypto requirements are unfulfilled.

\begin{lstlisting}[language=C, caption={NPNT Hardware Intercept Layer (\texttt{uav\_simulator.c})}]
// Inside process_command() validation matrix
if (cmd->command_id == KS_CMD_ARM) {
    if (g_npnt_enabled && !g_npnt_validated) {
        // NPNT Flight Clearance Hardware Locked. Rejecting ARM request.
        ack.result = KS_ACK_REJECTED; 
        return;
    }
    // Proceed to standard engine cryptographic arming ...
}
\end{lstlisting}

\subsubsection{Intended Protocol Behavior and Impact}
The UAV boots linearly into a locked state. The GCS must transmit the encrypted \texttt{KS\_MSG\_NPNT\_PA}. The UAV internally decrypts and validates the Ed25519 signature inside the payload against the known DGCA pub-key. Upon success, the system digitally unlocks.

\vspace{0.5cm}

\subsection{ASTM F3411: Remote ID Explicit Beacons}

\subsubsection{What It States}
Mandates an \textbf{Open Broadcast} concept. UAVs must physically unmask their serialized identifier vectors continuously over public bands (WiFi/Bluetooth) for law enforcement observability.

\subsubsection{Relevance to Kestrel}
This created a paradoxical conflict with Kestrel's core mandate: absolute zero-trust packet concealment. Leaking data could theoretically compromise military C2 tunnels.

\subsubsection{Implemented Code Changes}

We isolated the module execution completely outside the session tracker and hard-coded mathematical demotions into the base policy allocator.

\begin{lstlisting}[language=C, caption={Strict Cryptographic Policy Demotion (\texttt{kestrel.c})}]
// Inside kestrel.c encryption allocation paths
ks_encrypt_policy_t ks_get_encrypt_policy(uint16_t msg_id) {
    switch (msg_id) {
        case KS_MSG_RID_BASIC_ID:
        case KS_MSG_RID_LOCATION:
            // CRITICAL: MUST demote cipher application for law enforcement
            return KS_ENCRYPT_NEVER; 
        case KS_MSG_NPNT_PA:
            return KS_ENCRYPT_ALWAYS;
        // ... switch drops to optional bounds (ChaCha20 defaults)
    }
}
\end{lstlisting}

\subsubsection{Intended Protocol Behavior and Impact}
Kestrel now functions as a dual-channel router. Standard communications operate inside an impenetrable Poly1305 crypt-shell. Simultaneously, F3411 location updates are squawked broadly in plaintext, providing 100\% operational alignment without disturbing C2 integrity.

\vspace{0.5cm}

\subsection{NATO STANAG 4586: UAV--GCS Interoperability}

\subsubsection{What It States}
Defines standard interfaces between Unmanned Aerial Vehicles and Common Universal Control Systems (CUCS), requiring specialized Data Link routing mechanisms and standardized identifiers for heterogeneous drone swarms.

\subsubsection{Relevance to Kestrel}
Kestrel historically operated on localized formats. To support Level-of-Interoperability (LOI) 2-5 scaling, we needed proper header routing mapping components and larger addressable bounds.

\subsubsection{Implemented Code Changes}
We structured the extended packet headers to support \texttt{sys\_id} and \texttt{target\_sys\_id} partitioning, establishing direct compliance scaffolding.

\begin{lstlisting}[language=C, caption={STANAG Interoperability Extensible Header Hooks (\texttt{kestrel.h})}]
// Extended Header
uint8_t sys_id;        // 6-bit (Source System Identifier)
uint8_t comp_id;       // 4-bit (VSM/Payload Component Mapping)
uint8_t target_sys_id; // 6-bit (0 = CUCS Broadcast)
uint16_t msg_id;       // 12-bit (Expanded to 4096 hooks)
\end{lstlisting}

\subsubsection{Intended Protocol Behavior and Impact}
Provides the fundamental routing foundation for future NATO-grade multi-agent swarming architectures. While full CUCS message matrices are incomplete ($\sim$30\% completion), the architectural hooks guarantee expansions without necessitating a rewrite of \texttt{kestrel\_fast.c}.

\vspace{0.5cm}

\subsection{NATO STANAG 4609: Digital Motion Imagery}

\subsubsection{What It States}
Mandates digital motion imagery (DMI) platforms to embed standard KLV (Key-Length-Value) telemetry metadata synchronized sequentially with streaming video.

\subsubsection{Relevance to Kestrel}
Though Kestrel operates beneath the video aggregation layer, it must maintain a completely isolated data stream to prevent high-bandwidth KLV metadata from choking critical C2 telemetry.

\subsubsection{Implemented Code Changes}
We implemented active KLV (Key-Length-Value) metadata packaging using the explicitly reserved \texttt{KS\_STREAM\_VIDEO} identifier. Real-time UAV telemetry is converted into ST0601 compliant KLV bytes and multiplexed synchronously inside an MPEG-TS container payload.

\begin{lstlisting}[language=C, caption={STANAG 4609 ST0601 KLV Multiplexing (\texttt{uav\_simulator.c})}]
// STANAG 4609 MPEG-TS Video
ks_klv_uav_state_t klv_state = {
    .timestamp_us = (uint64_t)time(NULL) * 1000000ULL,
    .mission_id   = "KESTREL-SORTIE-001",
    .heading_deg  = state.yaw,
    // ... [Attitude & Geo Variables]
};

uint8_t klv_buf[256];
int klv_len = ks_klv_build_st0601(klv_buf, sizeof(klv_buf), &klv_state);

// Provide KLV metadata PES packet multiplexed into TS
ts_len += ks_ts_mux_write_pes(&ts_mux, KS_TS_PID_KLV, 0xBD, klv_state.timestamp_us, klv_buf, klv_len, &ts_buf[ts_len], MAX_TS);

// Package inside KS_STREAM_VIDEO wrapper
ks_header_t hdr_ts = {0};
hdr_ts.stream_type = KS_STREAM_VIDEO;
hdr_ts.msg_id = KS_MSG_VIDEO_TS;
\end{lstlisting}

\subsubsection{Intended Protocol Behavior and Impact}
The architecture now sits at roughly $\sim$85\% completion for Digital Motion Imagery standards. The Kestrel runtime successfully packages precise spatial metadata into compliant KLV elements, multiplexes them into standard MPEG-TS transport packets alongside video stream descriptors (PAT/PMT), and successfully transmits them via a localized, completely isolated socket without flooding critical C2 telemetry paths.


% ============================================================
% 4. CONCLUSION
% ============================================================
\subsection{JARUS SORA OSO\#06: C2 Link Cyber Security}

\subsubsection{What It States}
The JARUS SORA (Specific Operations Risk Assessment) framework, Annex E --- \emph{Cyber Safety Extension} --- defines \textbf{Operational Safety Objective \#06}: every C2 link must provide Confidentiality, Integrity, and Mutual Authentication (CIA triad) at a robustness level commensurate with the operation's SAIL (Specific Assurance and Integrity Level). All SAIL III and above require verifiable, time-stamped security event logging that can be produced for regulatory audit.

\subsubsection{Relevance to Kestrel}
Kestrel's core already satisfies the CIA triad:
\begin{itemize}
    \item \textbf{Confidentiality:} ChaCha20-Poly1305 AEAD, 256-bit session key.
    \item \textbf{Integrity:} 128-bit Poly1305 MAC bound to header as Additional Authenticated Data.
    \item \textbf{Anti-replay:} 32-bit sliding window (64-bit in Legion mode).
    \item \textbf{Mutual Authentication:} X25519 ECDH key exchange plus Ed25519 identity signatures on \texttt{KS\_MSG\_KEY\_EXCHANGE}.
\end{itemize}
The sole gap was \textbf{security event logging} --- OSO\#06 requires that replay attacks, MAC failures, and authentication outcomes be recorded in a format producible for operational audit.

\subsubsection{Implemented Code Changes}

We created a pure additive shim (\texttt{kestrel\_sora.c} / \texttt{kestrel\_sora.h}) with zero modifications to any core protocol file.

\begin{lstlisting}[language=C, caption={SORA OSO\#06 Event Types (\texttt{kestrel\_sora.h})}]
typedef enum {
    KS_SORA_REPLAY_REJECTED  = 0x01, // Replay window rejection
    KS_SORA_MAC_FAIL         = 0x02, // Poly1305 authentication failure
    KS_SORA_MUTUAL_AUTH_OK   = 0x03, // ECDH + EdDSA handshake succeeded
    KS_SORA_MUTUAL_AUTH_FAIL = 0x04, // EdDSA signature verification failed
    KS_SORA_KEY_ROTATED      = 0x05, // Session key established
    KS_SORA_LINK_ANOMALY     = 0x06, // Unexplained link gap
} ks_sora_event_t;
\end{lstlisting}

A 64-entry circular ring buffer records each event (timestamp, sys\_id, sequence number, error code). The compliance assertion evaluates all OSO\#06 conditions in O(1):

\begin{lstlisting}[language=C, caption={OSO\#06 Compliance Gate (\texttt{kestrel\_sora.c})}]
bool ks_sora_is_compliant(const ks_sora_ctx_t *ctx) {
    if (ctx->auth_ok_count < KS_SORA_AUTH_REQUIRED)    return false;
    if (ctx->mac_fail_count >= KS_SORA_MAC_FAIL_THRESHOLD) return false;
    if (ctx->replay_count >= KS_SORA_REPLAY_THRESHOLD) return false;
    if (ctx->auth_fail_count > 0 &&
        ctx->auth_ok_count == 0)                       return false;
    return true;
}
\end{lstlisting}

Call-site hooks were added in \texttt{uav\_simulator.c} and \texttt{gcs\_receiver.c} to tap into existing return paths without altering any logic:

\begin{lstlisting}[language=C, caption={Parse Error Hook (both simulators)}]
// After every ks_parse_char_zerocopy() call
ks_sora_on_parse_result(&g_sora_ctx, result,
                         sys_id, seq, get_time_ms());
// After EdDSA signature fails
ks_sora_log(&g_sora_ctx, KS_SORA_MUTUAL_AUTH_FAIL,
             get_time_ms(), sys_id, seq, 1);
// After KS_ECDH_ESTABLISHED is set
ks_sora_log(&g_sora_ctx, KS_SORA_MUTUAL_AUTH_OK,
             get_time_ms(), sys_id, seq, 0);
\end{lstlisting}

\subsubsection{CI/CD Gate and Unit Tests}
A 10-case unit test suite (\texttt{testing/sora\_unit\_test.c}) exercises all 43 compliance assertions including ring-wrap behaviour, threshold boundaries, and all three compliance state transitions. The test runs on every GitHub Actions push under the \texttt{build-linux} job with \texttt{-Wall -Werror}, producing a hard build failure on any regression.

\subsubsection{Intended Protocol Behavior and Impact}
Every session now produces a verifiable, time-stamped security audit trail. The \texttt{ks\_sora\_is\_compliant()} API provides an in-flight compliance gate that operators can assert before authorizing BVLOS operations. The ring buffer write path is a single struct copy plus counter increment --- zero measurable latency overhead on the C2 loop.

\vspace{0.5cm}

\section{Conclusion}

Through meticulous additive-shim engineering, Kestrel transforms effortlessly from a highly insulated experimental cryptographic transport protocol into a rigorously verified, multi-compliant aerospace framework. The codebase preserves its defining characteristics---zero-copy speed and resilient encryption pools---while safely executing multi-layered code compliance checks across six distinct regulatory domains: RTCA DO-362A, DO-377A, DGCA NPNT, ASTM F3411, NATO STANAG 4609, and JARUS SORA OSO\#06.

However, it is structurally imperative to recognize that the schematics and compliance metrics outlined herein are exclusively reference models and do not constitute a formally certified source-code audit. Native translation across heterogeneous flight platforms inherently introduces unverified code variations, latency bugs, and integration divergences outside these theoretical baselines, which directly impact absolute compliance. Consequently, formal, independent algorithmic scrutinization and rigorous code-path auditing are mandatory prior to any operational deployment. The reference architecture presented inherently disclaims all direct, indirect, or implied liability, and the authors assume no technical or organizational responsibility should runtime implementations develop defects or fail to satisfy empirical aviation compliance thresholds.

\end{document}
